\begin{table}[H]
\centering
\scriptsize
\setlength{\tabcolsep}{3pt}
\begin{tabular}{l|ccc|ccc|ccc}
\toprule
 & \multicolumn{3}{c}{IN} & \multicolumn{3}{c}{C100} & \multicolumn{3}{c}{C10} \\
model & Gauss. & Haar & PC-perm. & Gauss. & Haar & PC-perm. & Gauss. & Haar & PC-perm. \\
\midrule
ViT-T & $-0.005$ & $-0.014$ & $-0.038$ & $-0.044$ & $-0.031$ & $-0.055$ & $-0.057$ & $-0.027$ & $-0.050$ \\
ViT-S & $-0.023$ & $-0.025$ & $-0.052$ & $-0.034$ & $-0.025$ & $-0.049$ & $-0.098$ & $-0.063$ & $-0.090$ \\
ViT-B & $-0.027$ & $-0.025$ & $-0.054$ & $-0.018$ & $-0.011$ & $-0.039$ & $-0.090$ & $-0.054$ & $-0.078$ \\
ViT-L & $-0.028$ & $-0.029$ & $-0.040$ & $-0.047$ & $-0.034$ & $-0.051$ & $-0.104$ & $-0.066$ & $-0.096$ \\
\midrule
DINO-B & $-0.005$ & $-0.002$ & $-0.008$ & $-0.068$ & $-0.056$ & $-0.068$ & $-0.100$ & $-0.070$ & $-0.097$ \\
DINOv2-S & $+0.010$ & $+0.015$ & $+0.013$ & $-0.050$ & $-0.029$ & $-0.051$ & $-0.093$ & $-0.083$ & $-0.098$ \\
DINOv2-B & $+0.004$ & $+0.015$ & $+0.006$ & $-0.042$ & $-0.007$ & $-0.041$ & $-0.132$ & $-0.120$ & $-0.124$ \\
DINOv2-L & $-0.013$ & $+0.000$ & $-0.011$ & $-0.053$ & $-0.008$ & $-0.044$ & $-0.153$ & $-0.136$ & $-0.122$ \\
DINOv2-G & $-0.000$ & $+0.016$ & $+0.001$ & $-0.080$ & $-0.032$ & $-0.074$ & $-0.155$ & $-0.138$ & $-0.110$ \\
\midrule
CLIP-B & $-0.028$ & $-0.028$ & $-0.048$ & $-0.046$ & $-0.036$ & $-0.053$ & $-0.079$ & $-0.045$ & $-0.075$ \\
CLIP-L & $-0.016$ & $-0.017$ & $-0.029$ & $-0.042$ & $-0.031$ & $-0.052$ & $-0.087$ & $-0.049$ & $-0.082$ \\
SigLIP-B & $-0.019$ & $-0.020$ & $-0.044$ & $-0.042$ & $-0.032$ & $-0.046$ & $-0.096$ & $-0.056$ & $-0.081$ \\
\bottomrule
\end{tabular}
\caption{Excess under three constructions of the spectrum-matched null (10 replicates each): Gaussian coefficients (the paper's), Haar-rotated coefficients (exact sample spectrum) and PC-permutation (exact per-component marginals). The Gaussian and PC-permutation constructions agree; the exact-sample-spectrum construction yields smaller excesses for low-effective-rank clouds on the small-$C$ sets (most visibly DINOv2 on CIFAR-100), so the size of those excesses is partly a property of the null hypothesis, while signs and the ImageNet readings are stable.}
\label{tab:b27-nullvariants}
\end{table}
